
\begin{abstract}{Tutorial – 1: From Prototype to Product Readiness: IndustrialDesign Considerations for Dual Active BridgeConverters}{%
		Dr. Shan, Dr. Aravind, Dr. Gourahari}{%
		Schneider Electric, Bangalore
	}{%
		\Tutorialtag}
	\textbf{Abstract}:
	The Dual Active Bridge (DAB) converter is a promising DC-DC converter topology for medium-high power applications, offering bidirectional power transfer, galvanic isolation, and high efficiency. While academic research emphasizes on novel modulation strategies, robust control methods, and efficiency improvement; the transition from a prototype to an industrial product requires a broader and more holistic perspective.
	
	In addition to core power electronics, this tutorial presents the design of a Three-Phase Dual Active Bridge (3ϕ-DAB) converter for industrial applications, emphasizing the broader requirements that define product readiness. Critical aspects such as thermal management, mechanical robustness, and electromagnetic compliance (EMC/EMI) significantly influence industrial viability. Equally important are Design for Manufacturing (DFM), Design for Assembly (DFA), and Design for Reliability (DFR), which ensure scalability, cost-effectiveness, and consistency in mass production. Reliability practices, including DMA, DFMEA, further strengthen product robustness and long-term performance. With rising labour costs, industries are increasingly dependent on robotic and automated assembly lines. Therefore, converter design must also be automation-friendly, supporting repeatable and reliable processes such as automated soldering, reflow compatibility. Considering all these factors, the converter design must ultimately adhere to competitive market requirements by achieving high efficiency while minimizing cost, weight, and volume.
	This tutorial intends to give a brief insight to all the above said factors taking the Three Phase DAB as an example.
\end{abstract}

\begin{abstract}{Tutorial 2: Wireless Charging of Electric Vehicles: Design and Implementation}{%
		Dr. Francesca, Dr. Calina}{%
		Assistant Professor, Eindhoven University of Technology (TU/e), The Netherlands}{%
		\Tutorialtag}
	\textbf{Abstract}:
	Let’s assume that you want to design and implement a wireless charging system based on initial specifications and requirements. Where do we start the process? How do we choose the coil topology, compensation network, and power conversion stages? How do you model tolerances and evaluate/predict performance ? 
	At the end of this tutorial lecture, the student will be able to:
	Define boundary conditions and make design choices for the subcomponents of the wireless charger.
	Optimize the system's parameters.
	Evaluate the design considering manufacturing tolerances and other non-idealities. 
	
	Delivery methods:
	Lecture with slides
	Guided design task
\end{abstract}

%Tutorial – 3: 
\begin{abstract}{Tutorial 3: Reliability and circularity in the future highly efficient and high-power density DC power electronic converters}{%
		Dr. Riccardo\textsuperscript{1}, Dr. Mattia\textsuperscript{1}, Dr. Husev\textsuperscript{2}}{%
		University of Bologna, Italy\textsuperscript{1}, Gdansk University of Technology, Poland\textsuperscript{2}}{%
		\Tutorialtag}
	\textbf{Abstract}:
	As the global energy landscape evolves toward decarbonization, electrification, and digitalization, the demand for highly efficient and compact DC power electronic converters continues to surge. These converters are pivotal in enabling next-generation applications across renewable energy systems, electric mobility, data centers, and aerospace technologies. However, achieving high power density and efficiency often brings challenges related to thermal management, component stress, and long-term system reliability.
	
	This keynote will explore the interplay between reliability and circularity in the design and operation of future power converters. It will examine recent advances in wide-bandgap semiconductors, intelligent control strategies, and modular topologies that address efficiency and density goals while also focusing on design for longevity, ease of maintenance, reuse, and lifecycle carbon footprint.
	
	Bridging academic insights with industrial relevance, the keynote will outline strategies for embedding circular economy principles into the converter lifecycle—spanning design, usage, and lifecycle assessment. Attendees will gain a forward-looking perspective on how engineering innovation and sustainable design practices can together pave the way for resilient, resource-efficient power electronics in a rapidly transforming energy ecosystem.
\end{abstract}


\begin{abstract}{Tutorial 4: SiC MOSFETs Based Power Electronics: Challengesand Opportunities}{%
		Prof. Sandeep Anand\textsubscript{1}, Dr. Abhinav\textsubscript{2}}{%
	IIT Bombay\textsubscript{1} and IIT Bhubaneswar\textsubscript{2}, India
}{%
		\Tutorialtag}
	\textbf{Abstract}:
	The transition from conventional silicon IGBTs to silicon carbide (SiC) MOSFETs
	is redefining modern power electronics by enabling faster switching and higher switching
	frequency operation. This tutorial will explore the fundamental need for higher switching
	frequencies to achieve reduced passive component size while considering overall system
	efficiency. We will delve into the I-V characteristics of SiC MOSFETs, their switching
	behaviour, and the impact of parasitic capacitances on switching losses. A detailed comparison
	between SiC MOSFETs and Si-IGBTs will highlight the clear performance advantages of
	wide-bandgap semiconductors, while also addressing critical design challenges.
	Attendees will gain insights into circuit-level issues such as dv/dt-induced false turn-on,
	EMI/EMC concerns, gate driver complexities, and the need for fast-acting short-circuit
	protection. We will also discuss key reliability concerns, including gate oxide and body diode
	degradation mechanisms. The tutorial will also include the medium-voltage applications of the
	SiC MOSFETs, where these challenges become even more pronounced. The session will
	conclude with a brief presentation of our recent research findings, offering practical
	perspectives and design recommendations. This tutorial aims to equip engineers and
	researchers with a clear understanding of both the opportunities and obstacles in harnessing the
	full potential of SiC MOSFETs in next-generation power electronic systems.
\end{abstract}